\documentclass[10pt,a4paper]{amsart}
\usepackage[margin=19mm]{geometry}
\usepackage{amsmath,amssymb,mathtools}
\usepackage[scr=rsfs,cal=euler]{mathalfa}
\usepackage{xfrac}
\usepackage[unicode,hidelinks,bookmarks=false]{hyperref}
\newcommand{\ie}{i.e.\ }
\newcommand{\cf}{cf.\ }
\newcommand{\resp}{respectively\ }
\newcommand{\Subsection}{\subsection}
\providecommand{\nobreakdash}{\nobreak\hbox{-}}
\setlength{\emergencystretch}{2em}
\multlinegap0pt
\usepackage{mathtools}
\usepackage{amssymb}
\usepackage{algorithm}\usepackage{algpseudocode}
\usepackage{tikz}
\usepackage[shortlabels]{enumitem}
\usepackage[normalem]{ulem}
\usepackage{xcolor}
\usepackage{float}
\newtheorem{lemma}{Lemma}
\newtheorem{theorem}{Theorem}
\def\A{\mathcal{A}}
\def\B{\mathcal{B}}
\def\D{\mathcal{D}}
\DeclareMathOperator{\Diag}{Diag}
\def\K{\mathcal{K}}
\def\R{\mathbb{R}}
\def\Range{{\operatorname{Range}}}
\def\S{\mathbb{S}}
\def\Z{\mathcal{Z}}
\def\bQ{\overline{Q}}
\DeclareMathOperator{\diag}{diag}
\title{Sparsity-Cone SDP Relaxations and Applications to Variable Fixing for Sparse Quadratic Programs}
\author{Di Hou, Thai P.D. Nguyen, Kim-Chuan Toh, Guanyi Wang}
\date{Source: arXiv:2606.22894v1 (CC BY 4.0)}
\begin{document}
\maketitle

\noindent\emph{Source and licence.} Sparsity-Cone SDP Relaxations and Applications to Variable Fixing for Sparse Quadratic Programs. Version arXiv:2606.22894v1 (CC BY 4.0), distributed by arXiv under the Creative Commons Attribution 4.0 International licence, http://creativecommons.org/licenses/by/4.0/, as recorded in the arXiv licence metadata for this exact version and shown on the abstract page https://arxiv.org/abs/2606.22894v1. Full licence notice: \texttt{licenses/arxiv-2606.22894-CC-BY-4.0.txt}.\par\smallskip

\noindent\textit{Editorial scope.} Counted unit: the article's theorem thm:equivalence (source0.tex lines 1095--1108). The article's complete original proof of it is delivered with it (source0.tex lines 1110--1210, one \texttt{proof} environment of 101 lines). No mathematics is written, repaired or re-derived: the statement and the proof are the article's own lines, in the article's own notation, and no retained line is substituted or re-flowed.  Retained uncounted as the source-local context the proof needs: lines 243--253; lines 265--275; lines 339--347; lines 560--574; lines 969--982; lines 989--1002; lines 1028--1040.  Nothing was omitted from any retained span, so the package carries no omission record.  No retained line contains a citation, so the wrapper carries no bibliography and no bibliography entry.  No retained line contains a figure environment, an image-inclusion command or drawing code, so no image is delivered and the package declares no asset dependency.  Editorial formatting only, in addition: an AMS \texttt{amsart} preamble that restates the article's own declarations and macros used by the retained lines, namely the environments \texttt{lemma}, \texttt{theorem} on separate counters, together with the standard packages those lines need.  The retained environments print in this document as Lemma 1, Theorem 1 in order of appearance, whereas the article prints the same items with the numbering of its own full document.  The complete native source of the article as distributed in the arXiv e-print for this exact version is bundled unchanged as \texttt{sources/203390/source0.tex}.
\begin{equation}\label{eq:QP}
v^{\mathrm{P}}
\;\coloneqq\;
\min_{x}
\left\{
 x^\top Q x + 2c^\top x
\;\middle|\;
Ax=b,\; Bx\geq d,\;
\|x\|_0\le k
\right\},\tag{P}
\end{equation}

\begin{equation}\label{eq:MIQP}
\min_{x,z}
\left\{
 x^\top Q x + 2c^\top x
\;\middle|\;
Ax=b,\; Bx\geq d,\;
x\circ (1-z)=0,\;
e^\top z\leq k,\;
z\in \{0,1\}^n
\right\}.
\end{equation}

\begin{equation}
\label{eq:intro-Decomposed-SDP}
\min_{Y}\left\{
\langle \bQ,Y \rangle  \;\middle|\;  
 \A(Y)=0,\; \B(Y)\geq 0,\; 
Y_{11}=1,\;
 Y\in \K\cap \S^{n+1}_+
\right\},\tag{Q}
\end{equation}

\begin{lemma}
\label{lemma:equivalence}
The following equivalence holds for any $Y=[1,x^\top; x, X]\in\S^{n+1}_+$:
\begin{equation}\label{eq:kDiag-xx-general}
   Y\in\K
   \;\; \Longleftrightarrow\;\;
    \D(Y)\succeq 0
\;\; \Longleftrightarrow\;\;
k\Diag(X)\succeq xx^\top 
\;\; \Longleftrightarrow\;\;
    \sum_{i=1}^n \frac{x_i^2}{X_{ii}}\leq k.
\end{equation}
Note that for a nonnegative scalar $\alpha$, we define 
$\alpha/0 = 0$ if $\alpha = 0$, and $\alpha/0 = +\infty$ if $\alpha > 0$.
\end{lemma}

\begin{equation}\label{eq:Shor}
\min\left\{
\langle Q,X_{xx} \rangle+2c^\top x  \;\middle|\;
\begin{aligned}
&Ax=b,\;Bx\geq d,\; x = \diag(X_{xz}),\; z = \diag(X_{zz})\\
& 
\begin{bmatrix}
    1&x^\top & z^\top\\
    x&X_{xx} & X_{xz}\\
    z&X_{xz}^\top & X_{zz}
\end{bmatrix}\succeq 0,\; z\in \Z
\end{aligned}
\right\},
\end{equation}

\begin{equation}\label{eq:Shor-RLT}
\min\left\{
\langle Q,X_{xx} \rangle+2c^\top x  \;\middle|\;
\begin{aligned}
&\A(Y)=0,\; \B(Y)\geq 0,\; x = \diag(X_{xz}),\; z = \diag(X_{zz})\\
& 
\begin{bmatrix}
    1&x^\top & z^\top\\
    x&X_{xx} & X_{xz}\\
    z&X_{xz}^\top & X_{zz}
\end{bmatrix}\succeq 0,\; z\in \Z
\end{aligned}
\right\}.
\end{equation}

\begin{equation}
\label{eq:Decomposed-SDP-z}
\min_{Y,z}\left\{
\langle \bQ,Y \rangle  \;\middle|\;  
\begin{aligned}
&\A(Y)=0,\; \B(Y)\geq 0,\;Y_{11}=1,\; Y\succeq 0\\
& \begin{bmatrix}
z_i & x_i\\
x_i & X_{ii}
\end{bmatrix}\succeq 0,\ i\in [n],\;z\in \mathcal{Z}
\end{aligned}
\right\},
\end{equation}

\begin{theorem}
\label{thm:equivalence} 
The following inequality chain holds:
\[
\emph{val} \eqref{eq:Shor}
\;\le\;
\emph{val} \eqref{eq:Shor-RLT}
\;=\;
\emph{val} \eqref{eq:Decomposed-SDP-z}
\;=\;
\emph{val} \eqref{eq:intro-Decomposed-SDP}
\;\le\; \emph{val} \eqref{eq:QP}.
\]
\end{theorem}

\begin{proof}\;
We first compare the SDP--RLT relaxation \eqref{eq:Shor-RLT} with the reduced SDP--RLT relaxation \eqref{eq:Decomposed-SDP-z}. Let a feasible solution of \eqref{eq:Shor-RLT} be given and set $Y=[1, x^\top; x, X_{xx}]$.
Since $Y$ is a principal submatrix of the lifted SDP matrix, $Y\succeq0$. Moreover, each principal submatrix indexed by $(x_i,z_i)$ gives
\[
\begin{bmatrix}
X_{xx,ii}&X_{xz,ii}\\
X_{xz,ii}&X_{zz,ii}
\end{bmatrix}
=
\begin{bmatrix}
X_{xx,ii}&x_i\\
x_i&z_i
\end{bmatrix}
\succeq0,
\]
where $X_{xx,ii}$, $X_{xz,ii}$, and $X_{zz,ii}$ are the $(i,i)$-entries of $X_{xx}$, $X_{xz}$, and $X_{zz}$, respectively.
After permuting the two rows and columns, this is exactly the perspective LMI in \eqref{eq:Decomposed-SDP-z}. The RLT constraints depend only on $Y$, and the objective values agree because $\langle Q,X_{xx}\rangle+2c^\top x=\langle \bQ,Y\rangle$.
Hence every feasible solution of \eqref{eq:Shor-RLT} projects to a feasible solution of \eqref{eq:Decomposed-SDP-z} with the same objective value.

Conversely, let $(Y,z)$ be feasible for \eqref{eq:Decomposed-SDP-z}, where $Y=[1,x^\top;x,X]\succeq0$. Take a square factorization $Y=PP^\top$ with $P\in\R^{(n+1)\times(n+1)}$, and write the rows of $P$ as $p_0,\ldots,p_n\in\R^{n+1}$. For each $i\in[n]$, consider
\[
G_i=
\begin{bmatrix}
1&x_i&z_i\\
x_i&X_{ii}&x_i\\
z_i&x_i&z_i
\end{bmatrix}.
\]
The diagonal entries and the principal $2\times2$ minors of $G_i$ are nonnegative because $Y\succeq0$, $z_i\in[0,1]$, and $[z_i,x_i;x_i,X_{ii}]\succeq0$. Its determinant is $\det(G_i)=(1-z_i)(z_iX_{ii}-x_i^2)\ge0$. Since a symmetric $3\times3$ matrix is positive semidefinite if and only if all of its principal minors are nonnegative, we have $G_i\succeq0$.
To construct the required new row explicitly, let $P_i$ be the two-row matrix with rows $p_0$ and $p_i$, and set
\[
S_i:=P_iP_i^\top=[1,x_i;x_i,X_{ii}],
\qquad
b_i:=(z_i,x_i)^\top .
\]
Since $G_i=[S_i,b_i;b_i^\top,z_i]\succeq0$, the generalized Schur complement gives $b_i\in\Range(S_i)$ and $z_i-b_i^\top S_i^\dagger b_i\ge0$. Define $u_i=P_i^\top S_i^\dagger b_i\in\R^{n+1}$. Then
\[
P_i u_i=S_iS_i^\dagger b_i=b_i,
\qquad
\|u_i\|^2=b_i^\top S_i^\dagger b_i .
\]
Let $\delta_i:=z_i-\|u_i\|^2\ge0$. Embed the factor rows in $\R^{2n+1}$ by replacing each $p_j$ with $(p_j,0)$, and continue to write $P$ for the matrix formed by these zero-padded rows. If $e_i$ is the $i$-th coordinate vector in $\R^n$, set
\[
q_i:=(u_i,\sqrt{\delta_i}\,e_i)\in\R^{2n+1}.
\]
Then $\langle p_0,q_i\rangle=z_i$, $\langle p_i,q_i\rangle=x_i$, and $\langle q_i,q_i\rangle=z_i$. The remaining inner products involving the vectors $q_i$ are not prescribed by \eqref{eq:Decomposed-SDP-z}; define $(X_{xz})_{ji}:=\langle p_j,q_i\rangle$ and $(X_{zz})_{ij}:=\langle q_i,q_j\rangle$.
Let $P_z\coloneq[q_1,\ldots,q_n]^\top$. Then
\[
\begin{bmatrix}
1&x^\top&z^\top\\
x&X&X_{xz}\\
z&X_{xz}^\top&X_{zz}
\end{bmatrix}
=
\begin{bmatrix}
P\\
P_z
\end{bmatrix}
\begin{bmatrix}
P\\
P_z
\end{bmatrix}^{\top}
\succeq0.
\]
Moreover,
\[
\diag(X_{xz})=x,\qquad \diag(X_{zz})=z.
\]
The RLT constraints are unchanged because they only involve $Y$, and the objective value is also unchanged. Hence \eqref{eq:Shor-RLT} and \eqref{eq:Decomposed-SDP-z} have the same optimal value. 

Next, we compare \eqref{eq:Decomposed-SDP-z} with the SC-SDP relaxation \eqref{eq:intro-Decomposed-SDP}. The two formulations have the same objective function and the same constraints involving $\A(Y)$, $\B(Y)$, $Y_{11}=1$, and $Y\succeq0$. Let $(Y,z)$ be feasible for \eqref{eq:Decomposed-SDP-z}. The LMI $[z_i,x_i;x_i,X_{ii}]\succeq0$ gives $z_i\ge0$ and $z_iX_{ii}\ge x_i^2$. Hence, with the convention used in Lemma~\ref{lemma:equivalence},
\[
\sum_{i=1}^n \frac{x_i^2}{X_{ii}}
\le
\sum_{i=1}^n z_i
=e^\top z
\le k .
\]
By Lemma~\ref{lemma:equivalence}, this is equivalent to $Y\in\K$. Thus $Y$ is feasible for \eqref{eq:intro-Decomposed-SDP}.

Conversely, let $Y$ be feasible for \eqref{eq:intro-Decomposed-SDP}. By Lemma~\ref{lemma:equivalence}, $\sum_{i=1}^n x_i^2/X_{ii}\le k$.
Define
\[
z_i:=
\begin{cases}
x_i^2/X_{ii}, & X_{ii}>0,\\
0, & X_{ii}=0.
\end{cases}
\]
Since $Y\succeq0$, we have $X_{ii}\ge x_i^2$, and hence $0\le z_i\le1$. The preceding inequality gives $e^\top z\le k$, so $z\in\mathcal Z$. By construction, $z_iX_{ii}=x_i^2$ for every $i$, and therefore
\[
\begin{bmatrix}
z_i&x_i\\
x_i&X_{ii}
\end{bmatrix}
\succeq0,\qquad i\in[n].
\]
Hence $(Y,z)$ is feasible for \eqref{eq:Decomposed-SDP-z}. Therefore $\emph{val} \eqref{eq:intro-Decomposed-SDP}=\emph{val} \eqref{eq:Decomposed-SDP-z}$.

Moreover, the SDP--RLT relaxation is obtained by adding RLT constraints to the SDP relaxation, so $\emph{val} \eqref{eq:Shor}\le \emph{val} \eqref{eq:Shor-RLT}$. Finally, since \eqref{eq:MIQP} is an exact reformulation of \eqref{eq:QP}, both the SDP and SDP--RLT problems are relaxations of the original problem, and hence $\emph{val} \eqref{eq:Shor-RLT}\le \emph{val} \eqref{eq:QP}$.
\end{proof}

\end{document}
